\documentclass[a4paper]{article}
\title{Standard Code Library}
\author{AntiLeaf}
\date{}

% green= '#10BC72' # '#C3E88D'
% cyan= '#00A1D6' # '#89DDFF'

% \usepackage{zxjatype}
% \setjamainfont{ipaexm.ttf}

\usepackage{graphicx, amssymb, amsmath, textcomp, booktabs}
% \usepackage[libertine,vvarbb]{newtxmath}
\usepackage[scr=rsfso]{mathalfa}
%\usepackage[lining,semibold,type1]{libertine} % a bit lighter than Times--no osf in math
\usepackage[T1]{fontenc} % best for Western European languages
\usepackage{minted}
\usepackage{listings, color, setspace, titlesec, fancyhdr, mdframed, multicol}
\usepackage{fontspec}
\usepackage{ucharclasses}
\usepackage{xunicode, xltxtra}
\usepackage{pdfpages}
\usepackage{tocloft}
\usepackage{nameref}
\usepackage{verbatim}
\usepackage{relsize}
\usepackage[normalem]{ulem}

\usepackage{color,xcolor}

\setlength{\itemindent}{0em}
\setlength\parindent{0em}

\definecolor{light-gray}{gray}{0.9}    % 1.灰度
\definecolor{black}{gray}{0.0}

\XeTeXlinebreaklocale "zh"
\XeTeXlinebreakskip = 0pt plus 1pt

%configure space between the two columns
\setlength{\columnsep}{13pt}

%configure fonts
\setmonofont{Consolas}[Scale=0.8]
\newfontfamily\substitutefont{等线}[Scale=0.9]
\setTransitionsForChinese{\begingroup\substitutefont}{\endgroup}

%configure minted to display codes
\definecolor{Gray}{rgb}{0.9,0.9,0.9}

%configure section style of table of content
\renewcommand\cftsecfont{\Large}

%configure section style
\titleformat{\section}
{\huge}			% The style of the section title
{\thesection.}				% a prefix
{4pt}						% How much space exists between the prefix and the title
{}					% How the section is represented
% \titleformat{\section}{\huge}{}{0pt}{}
% \titlespacing{\section}{0pt}{0pt}{0pt}
% \titlespacing{\subsection}{0pt}{0pt}{0pt}
% \titlespacing{\subsubsection}{0pt}{0pt}{0pt}

\usepackage{fancyhdr}
\usepackage[inner=1.35cm, outer=0.9cm, top=1.7cm, bottom=0.4cm]{geometry}

\pagestyle{fancy}
\setlength{\headsep}{0.1cm}
\setlength{\footskip}{0.7cm}

\chead{Standard Code Library}
\cfoot{\thepage}

\renewcommand{\headrulewidth}{0.5pt}
\renewcommand{\footrulewidth}{0.5pt}

\setminted[cpp]{
	style=materiallight,
	mathescape,
	linenos,
	autogobble,
	baselinestretch=0.9,
	tabsize=4,
	fontsize=\normalsize,
	%bgcolor=Gray,
	frame=single,
	framesep=1mm,
	framerule=0.3pt,
	numbersep=1mm,
	breaklines=true,
	breaksymbolsepleft=2pt,
	%breaksymbolleft=\raisebox{0.8ex}{ \small\reflectbox{\carriagereturn}}, %not moe!
	%breaksymbolright=\small\carriagereturn,
	breakbytoken=false,
	showtabs=true,
	tab={\relscale{1.08} $\color{light-gray}{\vert} \ \ \ $ \relscale{1}},
}

\setminted[python]{
	style=material,
	mathescape,
	linenos,
	autogobble,
	baselinestretch=0.9,
	tabsize=4,
	fontsize=\normalsize,
	%bgcolor=Gray,
	frame=single,
	framesep=0.8mm,
	framerule=0.3pt,
	numbersep=0.8mm,
	breaklines=true,
	breaksymbolsepleft=2pt,
	%breaksymbolleft=\raisebox{0.8ex}{ \small\reflectbox{\carriagereturn}}, %not moe!
	%breaksymbolright=\small\carriagereturn,
	breakbytoken=false,
	showtabs=true,
	tab={\relscale{1.08} $\color{light-gray}{\vert} \ \ \ $ \relscale{1}},
}

\begin{document}

	\begin{titlepage}
		
		\includepdf{cover.pdf}

	\end{titlepage}

	\begin{multicols}{2}
		
		\begin{spacing}{1}
			\tableofcontents
		\end{spacing}
		\newpage

	\end{multicols}

	\begin{multicols}{2}

		\section{数学}

			\subsection{插值}
				\subsubsection{牛顿插值}
					牛顿插值的原理是 \textbf{二项式反演}。

\paragraph{二项式反演}

$$ f(n) = \sum_{k = 0} ^ n {n \choose k} g(k) \; \iff \; g(n) = \sum_{k = 0} ^ n \left( -1 \right) ^ {n - k} {n \choose k} f(k) $$

可以用 $e^x$ 和 $e^{-x}$ 的麦克劳林展开式证明。

套用二项式反演的结论即可得到牛顿插值：

$$ f(n) = \sum_{i = 0} ^ k {n \choose i} r_i , \; \text{where} \; r_i = \sum_{j = 0} ^ i (-1) ^ {i - j} {i \choose j} f(j) $$

其中 $k$ 表示 $f(n)$ 的最高次项系数。

实现时右边的式子等价于 $k$ 次差分：

\inputminted{cpp}{../src/math/牛顿插值.cpp}

注意到预处理 $r_i$ 的式子满足卷积形式，必要时可以用 FFT 优化至 $O(k\log k)$。


				\subsubsection{拉格朗日插值}
					$$ f(x) = \sum_i f(x_i) \prod_{j \neq i} \frac {x - x_j} {x_i - x_j} $$

快速插值参见 \detailedref{PolyFastInterpolation}。

			
			\subsection{多项式}
				\subsubsection{FFT}
					\inputminted{cpp}{../src/math/FFT.cpp}

				\subsubsection{NTT}
					\inputminted{cpp}{../src/math/NTT.cpp}

				\subsubsection{任意模数卷积(三模数NTT)}
					\inputminted{cpp}{../src/math/任意模数卷积.cpp}

				% \subsection{毛梯梯}

			
				\subsubsection{多项式操作}
					\inputminted{cpp}{../src/math/多项式操作.cpp}

				\subsubsection{更优秀的多项式多点求值}
					这个做法不需要写求逆和取模, 但是神乎其技, 完全搞不懂原理 \\
					清空和复制之类的地方容易抄错, 抄的时候要注意
					\inputminted{cpp}{../src/math/更优秀的多项式多点求值.cpp}

				% \subsubsection{多项式快速插值}


				% \subsubsection{多项式复合与复合逆}


				\subsubsection{拉格朗日反演}
					用于求复合逆. \\
如果$f(x)$与$g(x)$互为复合逆，则有 \\
\begin{math}
\begin{aligned}[x^n]g(x)=\frac{1}{n}[x^{n-1}]\left(\frac{x}{f(x)}\right)^n\end{aligned} \\
\begin{aligned}[x^n]h(g(x))=\frac{1}{n}[x^{n-1}]h'(x)\left(\frac{x}{f(x)}\right)^n\end{aligned}
\end{math}

				% \subsubsection{分治FFT}


				\subsubsection{半在线卷积}
					\inputminted{cpp}{../src/math/半在线卷积.cpp}
				
				\subsubsection{常系数齐次线性递推$O(k\log k\log n)$}
					如果只有一次这个操作可以像代码里一样加上一个只求一次逆的优化, 否则就乖乖每次做完整的除法和取模
					\inputminted{cpp}{../src/math/常系数齐次线性递推.cpp}
				

			\subsection{FWT快速沃尔什变换}
				\inputminted{cpp}{../src/math/FWT.cpp}

			\subsection{单纯形}
				\inputminted{cpp}{../src/math/单纯形.cpp}

				% \subsubsection{线性规划对偶原理}


			\subsection{线性代数}
				% \subsubsection{高斯消元}


				% \subsubsection{任意模数高斯消元}


				% \subsubsection{自由元搜索}


				\subsubsection{线性基}
					% \inputminted{cpp}{../src/math/线性基.cpp}


				\subsubsection{线性代数知识}
					行列式: $$\det A = \sum_{\sigma} \text{sgn}(\sigma) \prod_{i}a_{i, \sigma_i}$$

逆矩阵: $$B = A^{-1} \iff AB = 1$$

代数余子式: $$C_{i, j} = (-1)^{i + j} M_{i, j} = (-1) ^ {i + j} \left| A ^ {i, j} \right|$$

也就是$A$去掉一行一列之后的行列式.

伴随矩阵: $$A^{*} = C^T$$

即代数余子式矩阵的转置.

同时我们有 $$ A^{*} = |A| A^{-1}$$

特征多项式: $$P_A(x) = \det{(Ix - A)}$$

特征根: 特征多项式的所有$n$个根(可能有重根).

			% \subsection{自适应Simpson积分}


			% \subsection{常见结论}
				% \subsubsection{线性齐次递推求通项}
			
				
			\subsection{常见数列}
				\subsubsection{伯努利数}
					\begin{math}
    \begin{aligned}B(x)=\sum_{i\ge 0}\frac{B_i x^i}{i!}=\frac x{e^x-1}\end{aligned} \\
    \begin{aligned}B_n=[n=0]-\sum_{i=0}^{n-1}{n\choose i}\frac{B_i}{n-k+1}\end{aligned} \\
    \begin{aligned}\sum_{i=0}^n{n+1\choose i}B_i=0\end{aligned} \\
    \begin{aligned}S_n(m)=\sum_{i=0}^{m-1}i^n=\sum_{i=0}^n{n\choose i}B_{n-i}\frac{m^{i+1}}{i+1}\end{aligned}
\end{math}

自然数幂次和关于次数的EGF:
$$ \begin{aligned} F(x)=&\sum_{k=0}^\infty \frac{\sum_{i=0}^n i^k}{k!}x^k\\ =&\sum_{i=0}^n e^{ix}\\ =&\frac{e^{(n+1)x-1}}{e^x-1} \end{aligned} $$
				
				\subsubsection{分拆数}
					\inputminted{cpp}{../src/math/分拆数.cpp}


		\section{数论}

			\subsection{$O(n)$预处理逆元}
				\inputminted{cpp}{../src/numbertheory/O(n)求逆元.cpp}

			\subsection{杜教筛}
				\inputminted{cpp}{../src/numbertheory/杜教筛.cpp}
			
			\subsection{线性筛}
				\inputminted{cpp}{../src/numbertheory/扩展线性筛.cpp}

			% \subsection{min_25筛}


			\subsection{Miller-Rabin}
				\inputminted{cpp}{../src/numbertheory/Miller-Rabin.cpp}

			\subsection{Pollard's Rho}
				\inputminted{cpp}{../src/numbertheory/Pollard-Rho.cpp}
			
			% \subsection{组合数取模}


			%\subsection{扩展欧几里得}


		\section{图论}

			\subsection{最小生成树}
				\subsubsection{Boruvka算法}
					思想: 每次选择连接每个连通块的最小边, 把连通块缩起来. \\
每次连通块个数至少减半, 所以迭代$O(\log n)$次即可得到最小生成树.

一种比较简单的实现方法: 每次迭代遍历所有边, 用并查集维护连通性和每个连通块的最小边权.

应用: 最小异或生成树
				
				\subsubsection{动态最小生成树}
					\inputminted{cpp}{../src/graph/动态最小生成树.cpp}
				
				% \subsubsection{最小树形图(朱刘算法)}

				
				\subsubsection{Steiner Tree 斯坦纳树}
					\paragraph{问题} 一张图上有 $k$ 个关键点，求让关键点两两连通的最小生成树。

\paragraph{做法} 状压 DP，$f_{i,S}$ 表示以 $i$ 号点为树根，$i$ 与 $S$ 中的所有点连通的最小边权和。

转移有两种：

\begin{enumerate}
	\item 枚举子集：
	$$ f_{i, S} = \min_{T\subset S} \left\{f_{i, T} + f_{i, S\setminus T}\right\} $$
	\item 新加一条边: 
	$$ f_{i, S} = \min_{(i, j)\in E} \left\{f_{j, S} + w_{i, j}\right\} $$
\end{enumerate}

第一种直接枚举子集 DP 就行了，第二种可以用 SPFA 或者 Dijkstra 松弛（显然负边一开始全选就行了，所以只需要处理非负边）。

复杂度 $O(n 3^k + 2^k \text{SSSP}(n, m)))$，其中 $\text{SSSP}(n,m)$ 可以是 $nm$（SPFA）或者 $n^2+m$ 或者 $m\log n$（后两种都是 Dijkstra）。

\inputminted{cpp}{../src/graph/steiner.cpp}

			
			
			\subsection{最短路}
				% \subsubsection{Dijkstra}
					
				
				% \subsubsection{差分约束}


				\subsubsection{k短路}
					\inputminted{cpp}{../src/graph/k短路.cpp}

			\subsection{仙人掌}
				\subsubsection{仙人掌DP}
					\inputminted{cpp}{../src/graph/仙人掌DP.cpp}
			
			\subsection{二分图}
				% \subsubsection{匈牙利}


				% \subsubsection{Hopcroft-Karp}


				\subsubsection{KM二分图最大权匹配}
					\inputminted{cpp}{../src/graph/KM二分图最大权匹配.cpp}

				% \subsubsection{二分图原理}
			
			
			\subsection{一般图匹配}
				\subsubsection{高斯消元}
					\inputminted{cpp}{../src/graph/基于线性代数的一般图匹配.cpp}

				\subsubsection{带花树}
					\inputminted{cpp}{../src/graph/带花树.cpp}

				\subsubsection{带权带花树}
					(有一说一这玩意实在太难写了, 抄之前建议先想想算法是不是假的或者有SB做法)
					\inputminted{cpp}{../src/graph/带权带花树.cpp}
			
			% \subsection{支配树}
			

			% \subsection{2-SAT}


			\subsection{最大流}
				\subsubsection{Dinic}
					\inputminted{cpp}{../src/graph/Dinic.cpp}

				\subsubsection{ISAP}
					\inputminted{cpp}{../src/graph/ISAP.cpp}

				\subsubsection{HLPP最高标号预流推进}
					\inputminted{cpp}{../src/graph/HLPP.cpp}
			
			\subsection{费用流}
				\subsubsection{SPFA费用流}
					\inputminted{cpp}{../src/graph/SPFA费用流.cpp}

				% \subsubsection{Dijkstra费用流}


				% \subsubsection{zkw费用流}
			

			% \subsection{网络流原理}
				% \subsubsection{最小割}


				% \subsubsection{费用流}


				% \subsubsection{常见建图方法}


				% \subsubsection{例题}
			

			% \subsection{弦图}
				% \subsubsection{完美消除序列、最大势算法}


				% \subsubsection{性质}
			

			% \subsection{其他}
				% \subsubsection{Stoer-Wagner全局最小割}


				% \subsubsection{最小割树}


				% \subsubsection{最大团搜索}
			

		\section{数据结构}

			\subsection{线段树}
				\subsubsection{非递归线段树}
					% \sout{让fstqwq手撕}

\begin{itemize}
	\item 如果$M = 2^k$, 则只能维护$[1, M - 2]$范围
	\item 找叶子: $i$对应的叶子就是$i + M$
	\item 单点修改: 找到叶子然后向上跳
	\item 区间查询: 左右区间各扩展一位, 转换成开区间查询
	\inputminted{cpp}{../src/datastructure/非递归线段树单点修改.cpp}

\end{itemize}
	
区间修改要标记永久化,并且求区间和和求最值的代码不太一样. \\

\textbf{区间加, 区间求和}
\inputminted{cpp}{../src/datastructure/非递归线段树区间加区间求和.cpp}

\textbf{区间加, 区间求最大值}
\inputminted{cpp}{../src/datastructure/非递归线段树区间加区间求最大值.cpp}

				\subsubsection{主席树}
					这种东西能不能手撕啊
	
			\subsection{陈丹琦分治}
				\inputminted{cpp}{../src/datastructure/CDQ分治.cpp}
	
			% \subsection{整体二分}


			% \subsection{块状链表}
	
	
			\subsection{Treap}
				\inputminted{cpp}{../src/datastructure/Treap.cpp}
				
				% \subsubsection{无旋Treap}


				% \subsubsection{可持久化Treap}
	
	
			\subsection{Splay}
				(参见LCT,除了\mintinline{cpp}{splay()}需要传一个点表示最终它的父亲,其他写法都和LCT相同)
				
			\subsection{树分治}
				% \subsubsection{静态树分治}

				
				\subsubsection{动态树分治}
					\inputminted{cpp}{../src/datastructure/动态树分治.cpp}

				\subsubsection{紫荆花之恋}
					\inputminted{cpp}{../src/datastructure/紫荆花之恋.cpp}
	
			\subsection{LCT}
				\subsubsection{不换根(弹飞绵羊)}
					\inputminted{cpp}{../src/datastructure/LCT(不换根).cpp}
			
				\subsubsection{换根/维护生成树}
					\inputminted{cpp}{../src/datastructure/LCT(换根).cpp}

				\subsubsection{维护子树信息}
					\inputminted{cpp}{../src/datastructure/LCT维护子树信息.cpp}
					
				\subsubsection{模板题:动态QTREE4(询问树上相距最远点)}
						\inputminted{cpp}{../src/datastructure/动态QTREE4.cpp}
	
			% \subsection{K-D树}
	
	
			% \subsection{树分块}
	
	
			% \subsection{树上莫队}
			
	
			\subsection{虚树}
				\inputminted{cpp}{../src/datastructure/虚树.cpp}
	
	
			\subsection{长链剖分}
				\inputminted{cpp}{../src/datastructure/长链剖分.cpp}
	
			\subsection{梯子剖分}
				\inputminted{cpp}{../src/datastructure/梯子剖分.cpp}
					
			\subsection{左偏树}
				(参见k短路)
	
			\subsection{常见根号思路}
				
\noindent{
	{\large\bfseries{通用}}
	\begin{itemize}
		\item 出现次数大于$\sqrt n$的数不会超过$\sqrt n$个
		\item 对于带修改问题, 如果不方便分治或者二进制分组, 可以考虑对操作分块, 每次查询时暴力最后的$\sqrt n$个修改并更正答案
		\item {\bfseries 根号分治}: 如果分治时每个子问题需要$O(N)$(N是全局问题的大小)的时间, 而规模较小的子问题可以$O(n^2)$解决, 则可以使用根号分治
			\begin{itemize}
				\item 规模大于$\sqrt n$的子问题用$O(N)$的方法解决, 规模小于$\sqrt n$的子问题用$O(n^2)$暴力
				\item 规模大于$\sqrt n$的子问题最多只有$\sqrt n$个
				\item 规模不大于$\sqrt n$的子问题大小的平方和也必定不会超过$n\sqrt n$
			\end{itemize}
		\item 如果输入规模之和不大于$n$(例如给定多个小字符串与大字符串进行询问), 那么规模超过$\sqrt n$的问题最多只有$\sqrt n$个
	\end{itemize}
}

\noindent{
	{\large\bfseries{序列}}
	\begin{itemize}
		\item 某些维护序列的问题可以用分块/块状链表维护
		\item 对于静态区间询问问题, 如果可以快速将左/右端点移动一位, 可以考虑莫队
			\begin{itemize}
				\item 如果强制在线可以分块预处理, 但是一般空间需要$n\sqrt n$ 
					\begin{itemize}
						\item 例题: 询问区间中有几种数出现次数恰好为$k$, 强制在线
					\end{itemize}
				\item 如果带修改可以试着想一想带修莫队, 但是复杂度高达$n^{\frac 5 3}$
			\end{itemize}
		\item 线段树可以解决的问题也可以用分块来做到$O(1)$询问或是$O(1)$修改, 具体要看哪种操作更多
	\end{itemize}
}

\noindent{
	{\large\bfseries{树}}
	\begin{itemize}
		\item 与序列类似, 树上也有树分块和树上莫队
			\begin{itemize}
				\item 树上带修莫队很麻烦, 常数也大, 最好不要先考虑
				\item 树分块不要想当然
			\end{itemize}
		\item 树分治也可以套根号分治, 道理是一样的
	\end{itemize}
}

\noindent{
	{\large\bfseries{字符串}}
	\begin{itemize}
		\item 循环节长度大于$\sqrt n$的子串最多只有$O(n)$个, 如果是极长子串则只有$O(\sqrt n)$个
	\end{itemize}
}

\noindent {
	{\large\bfseries{关于莫队}}

	莫队是可以改造成只有插入和撤销(或者只有删除和撤销)的版本的.

	例如维护dfs序时就可以使用链表, 配合只有删除的莫队就可以做到$O(n\sqrt n)$.

	另外如果$n$和$q$不平衡, 块大小应该设为$\frac n {\sqrt q}$.
}

		\section{字符串}

			\subsection{AC自动机}
				\inputminted{cpp}{../src/string/AC自动机.cpp}

			\subsection{后缀数组}
				\subsubsection{SA-IS}
					\inputminted{cpp}{../src/string/sais.cpp}
			
				\subsubsection{SAMSA}
				 	\inputminted{cpp}{../src/string/SAMSA.cpp}

			% \subsection{后缀平衡树}


			\subsection{后缀自动机}
				(广义后缀自动机复杂度就是$O\left(n\left|\Sigma\right|\right)$, 也没法做到更低了)
				\inputminted{cpp}{../src/string/后缀自动机.cpp}

			\subsection{回文树}
				\inputminted{cpp}{../src/string/回文树.cpp}

				\subsubsection{广义回文树}
					(代码是梯子剖分的版本,压力不大的题目换成直接倍增就好了,常数只差不到一倍)
					\inputminted{cpp}{../src/string/广义回文树.cpp}


			% \subsection{序列自动机}


			\subsection{Manacher马拉车}
				\inputminted{cpp}{../src/string/manacher.cpp}

			\subsection{KMP}
				

				\subsubsection{ex-KMP}
					\inputminted{cpp}{../src/string/exKMP.cpp}
			
			% \subsection{All Substring LCS}


			% \subsection{字符串原理}
				

			\section{动态规划}
				\subsection{决策单调性$O(n\log n)$}
					\inputminted{cpp}{../src/DP/决策单调性.cpp}

			\section{Miscellaneous}
				\subsection{$O(1)$快速乘}
					\inputminted{cpp}{../src/misc/O(1)快速乘.cpp}
				
				\subsection{$O(n^2)$高精度}
					\inputminted{cpp}{../src/misc/高精度.cpp}

				\subsection{xorshift}
					\inputminted{cpp}{../src/misc/xorshift.cpp}
				
				\subsection{枚举子集}
					(注意这是$t\ne 0$的写法, 如果可以等于$0$需要在循环里手动\mintinline{cpp}{break})
\begin{minted}{cpp}
for (int t = s; t; (--t) &= s) {
    // do something
}
\end{minted}
					
				\subsection{STL}
					\subsubsection{vector}
	\begin{itemize}
		\item \mintinline{cpp}{vector(int nSize)}: 创建一个vector, 元素个数为nSize
		\item \mintinline{cpp}{vector(int nSize, const T &value)}: 创建一个vector, 元素个数为nSize, 且值均为value
		\item \mintinline{cpp}{vector(begin, end)}: 复制[begin, end)区间内另一个数组的元素到vector中
		\item \mintinline{cpp}{void assign(int n, const T &x)}: 设置向量中前n个元素的值为x
		\item \mintinline{cpp}{void assign(const_iterator first, const_iterator last)}: 向量中[first, last)中元素设置成当前向量元素
		\item \mintinline{cpp}{void emplace_back(Args&&... args)}: 自动构造并push\_back一个元素, 例如对一个存储pair的vector可以 \mintinline{cpp}{v.emplace_back(x, y)}
	\end{itemize}

\subsubsection{list}
	\begin{itemize}
		\item \mintinline{cpp}{assign()} 给list赋值 
		\item \mintinline{cpp}{back()} 返回最后一个元素 
		\item \mintinline{cpp}{begin()} 返回指向第一个元素的迭代器 
		\item \mintinline{cpp}{clear()} 删除所有元素 
		\item \mintinline{cpp}{empty()} 如果list是空的则返回true 
		\item \mintinline{cpp}{end()} 返回末尾的迭代器
		\item \mintinline{cpp}{erase()} 删除一个元素
		\item \mintinline{cpp}{front()} 返回第一个元素
		\item \mintinline{cpp}{insert()} 插入一个元素到list中
		\item \mintinline{cpp}{max_size()} 返回list能容纳的最大元素数量
		\item \mintinline{cpp}{merge()} 合并两个list
		\item \mintinline{cpp}{pop_back()} 删除最后一个元素
		\item \mintinline{cpp}{pop_front()} 删除第一个元素
		\item \mintinline{cpp}{push_back()} 在list的末尾添加一个元素
		\item \mintinline{cpp}{push_front()} 在list的头部添加一个元素
		\item \mintinline{cpp}{rbegin()} 返回指向第一个元素的逆向迭代器
		\item \mintinline{cpp}{remove()} 从list删除元素
		\item \mintinline{cpp}{remove_if()} 按指定条件删除元素
		\item \mintinline{cpp}{rend()} 指向list末尾的逆向迭代器
		\item \mintinline{cpp}{resize()} 改变list的大小
		\item \mintinline{cpp}{reverse()} 把list的元素倒转
		\item \mintinline{cpp}{size()} 返回list中的元素个数
		\item \mintinline{cpp}{sort()} 给list排序
		\item \mintinline{cpp}{splice()} 合并两个list
		\item \mintinline{cpp}{swap()} 交换两个list
		\item \mintinline{cpp}{unique()} 删除list中重复的元
	\end{itemize}

\subsubsection{unordered\_set/map}

\begin{itemize}
	\item \mintinline{cpp}{unordered_map<int, int, hash>}: 自定义哈希函数, 其中hash是一个带重载括号的类.
\end{itemize}

\subsubsection{自定义 Hash}
	\inputminted{cpp}{../src/misc/hash.cpp}

				\subsection{pb\_ds}


				\subsection{rope}
					

			\section{注意事项}
				\subsection{常见下毒手法}
					\noindent
\begin{itemize}
    \item 高精度高低位搞反了吗
    \item 线性筛抄对了吗
    \item 快速乘抄对了吗
    \item \mintinline{cpp}{i <= n, j <= m}
    \item sort比较函数是不是比了个寂寞
    \item 该取模的地方都取模了吗
    \item 边界情况(+1-1之类的)有没有想清楚
    \item \bfseries{特判是否有必要, 确定写对了吗}
\end{itemize}

				\subsection{场外相关}
					\noindent
\begin{itemize}
    \item 安顿好之后查一下附近的咖啡店,打印店,便利店之类的位置,以备不时之需
    \item 热身赛记得检查一下编译注意事项中的代码能否过编译,还有熟悉比赛场地,清楚洗手间在哪儿,测试打印机(如果可以)
    \item 比赛前至少要翻一遍板子,尤其要看原理与例题
    \item 比赛前一两天不要摸鱼,要早睡,有条件最好洗个澡;比赛当天不要起太晚,维持好的状态
    \item 赛前记得买咖啡,最好直接安排三人份,记得要咖啡因比较足的;如果主办方允许,就带些巧克力之类的高热量零食
    \item 入场之后记得检查机器,尤其要逐个检查键盘按键有没有坏的;如果可以的话,调一下gedit设置
    \item 开赛之前调整好心态,比赛而已,不必心急.
\end{itemize}

				\subsection{做题策略与心态调节}
					\noindent
\begin{itemize}
	\item 拿到题后立刻按照商量好的顺序读题,前半小时最好跳过题意太复杂的题(除非被过穿了)
	
	\item 签到题写完不要激动,稍微检查一下最可能的下毒点再交,避免无谓的罚时
		\begin{itemize}
			\item 一两行的那种傻逼题就算了
		\end{itemize}
		
	\item 读完题及时输出题意,一方面避免重复读题,一方面也可以让队友有一个初步印象,方便之后决定开题顺序
	
	\item 如果不能确定题意就不要贸然输出甚至上机, 尤其是签到题, 因为样例一般都很弱
	
	\item —个题如果卡了很久又有其他题可以写,那不妨先放掉写更容易的题,不要在一棵树上吊死
		\begin{itemize}
			\item 不要被—两道题搞得心态爆炸,一方面急也没有意义,一方面你很可能真的离AC就差一步
		\end{itemize}
		
	\item 榜是不会骗人的,一个题如果被不少人过了就说明这个题很可能并没有那么难;如果不是有十足的把握就不要轻易开没什么人交的题;另外不要忘记最后一小时会封榜
	
	\item 想不出题/找不出毒自然容易犯困,一定不要放任自己昏昏欲睡,最好去洗手间冷静一下,没有条件就站起来踱步
	
	\item 思考的时候不要挂机,一定要在草稿纸上画一画,最好说出声来最不容易断掉思路
	
	\item 出完算法一定要check一下样例和一些trivial的情况,不然容易写了半天发现写了个假算法
	
	\item 上机前有时间就提前给需要思考怎么写的地方打草稿,不要浪费机时
	
	\item 查毒时如果最难的地方反复check也没有问题,就从头到脚仔仔细细查一遍,不要放过任何细节,即使是并查集和sort这种东西也不能想当然
	
	\item 后半场如果时间不充裕就不要冒险开难题,除非真的无事可做
		\begin{itemize}
			\item 如果是没写过的东西也不要轻举妄动,在有其他好写的题的时候就等一会再说
		\end{itemize}
		
	\item 大多数时候都要听队长安排,虽然不一定最正确但可以保持组织性
	
	\item 最好注意一下影响,就算忍不住嘴臭也不要太大声
	
	\item 任何时候都不要着急,着急不能解决问题,不要当喆国王
	
	\item 输了游戏,还有人生;赢了游戏,还有人生.
\end{itemize}

	\end{multicols}

\end{document}